\documentclass[14pt]{extreport}

\usepackage[margin=1in]{geometry}
\usepackage[T1]{fontenc}
\usepackage[utf8]{inputenc}
\usepackage{hyperref}

\begin{document}

\begin{center}
\textbf{Experiment 8: Markov Decision Processes and Dynamic Programming}\\[0.8em]
Manoj Kumar Panda -- 20251602013\\
Meetkumar Kankotiya -- 20251602007\\
Shrey Parikh -- 20251602017\\[0.5em]
Group--10\\
CS663 -- Artificial Intelligence\\
Indian Institute of Information Technology Vadodara (IIITV)
\end{center}

\vspace{1em}

\noindent\textbf{Abstract:}\\
This experiment studies sequential decision making in stochastic environments
using Markov Decision Processes (MDPs) and dynamic programming techniques.
We first apply value iteration to a stochastic gridworld to compute optimal
state values under different reward structures. Then we model the Gbike
bicycle rental problem as a finite MDP and solve it using policy iteration.
The implementation demonstrates how value functions, transition dynamics and
rewards interact to produce optimal policies.

\vspace{1em}

\noindent\textbf{Introduction:}\\
In many AI problems an agent must choose actions not just for immediate
reward, but also for their long-term consequences. Markov Decision Processes
provide a mathematical framework in which states, actions, transition
probabilities and rewards are explicitly modelled. Dynamic programming
algorithms such as value iteration and policy iteration exploit the Bellman
equations to compute optimal policies when a model of the environment is
available. In this experiment we implement these ideas for two problems:
a $4 \times 3$ gridworld and the Gbike bicycle rental problem.

\vspace{1em}

\noindent\textbf{Problem Description:}\\
The work is divided into two parts:
\begin{itemize}
  \item \textbf{Stochastic Gridworld:} The agent moves on a $4 \times 3$
  grid with two terminal states ($+1$ and $-1$) and one wall. At each step
  the agent chooses an intended direction (up, down, left, right), but due
  to stochasticity the actual move is the intended direction with probability
  $0.8$ and a perpendicular move with probability $0.1$ each. Bumping into a
  wall or boundary leaves the agent in the same state. All non-terminal states
  receive a reward $r(s)$ per step, and terminal states provide their terminal
  rewards. The goal is to compute the value function corresponding to the
  optimal policy for different choices of $r(s)$.
  \item \textbf{Gbike Bicycle Rental:} This is an adaptation of the Jack's
  car rental problem. The manager controls two rental locations with a
  limited number of bicycles at each. Each day, customer rental requests and
  returns at both locations are modelled as Poisson random variables with
  given means. Renting a bike yields a fixed positive reward, while moving
  bikes overnight between locations incurs a cost. The state is the number
  of bikes at each location at the end of the day, and the action is the
  number of bikes moved from one location to the other overnight. The task
  is to find an optimal stationary policy that maximises long-term expected
  discounted reward. A modified version of the problem introduces one free
  bike moved from location 1 to 2 and an extra parking cost if more than
  ten bikes are kept at a location overnight.
\end{itemize}

\vspace{1em}

\noindent\textbf{Algorithm:}\\
For the gridworld, we use \textbf{value iteration}. For each state $s$, the
update is
\[
V_{k+1}(s) = \max_a \sum_{s'} P(s' \mid s,a)\,\bigl(R(s,a,s') + \gamma V_k(s')\bigr),
\]
where $P(s' \mid s,a)$ is obtained from the stochastic transition model and
$R(s,a,s')$ is the immediate reward. Iteration continues until the maximum
change in $V(s)$ across all states falls below a small threshold. The optimal
policy is then obtained by acting greedily with respect to the final value
function.

For the Gbike problem we use \textbf{policy iteration}. Starting from an
initial policy, we alternate between:
\begin{itemize}
  \item \emph{Policy evaluation:} solving for the value function $V^\pi(s)$
  under the current policy $\pi$ using iterative updates based on the Bellman
  expectation equation.
  \item \emph{Policy improvement:} for each state $s$, choosing a new action
  $a$ that maximises the expected return
  \[
  q^{\pi}(s,a) = \mathbb{E}\bigl[R + \gamma V^\pi(S') \mid S=s, A=a\bigr].
  \]
\end{itemize}
If the policy does not change during an improvement step, it is optimal.
The Gbike environment uses truncated Poisson distributions to enumerate
possible numbers of rental requests and returns at each location, and the
expected return is computed by summing over these possibilities.

\vspace{1em}

\noindent\textbf{Implementation:}\\
The complete Python implementation (gridworld value iteration, base Gbike
policy iteration and the modified Gbike variant) for this experiment is
available at the following URL:\\[0.3em]
\hspace*{1em}\url{https://github.com/mkankotiya/AI-CS659/blob/main/gbike.py}

\vspace{1em}

\noindent\textbf{Results:}\\
For the gridworld, different choices of the step reward $r(s)$ produce
different optimal value functions and policies. When $r(s)$ is moderately
negative, the agent prefers shorter paths to the positive terminal state and
avoids unnecessary wandering. As $r(s)$ becomes more negative, the agent
tries to terminate episodes quickly even if that increases the chance of
reaching the $-1$ terminal state. When $r(s)$ is positive, the agent may
prefer to cycle in the grid rather than terminate, since continuing to move
yields more cumulative positive reward.

In the Gbike base problem, policy iteration converges to a stable policy that
balances moving bicycles overnight against expected rental income. Typically,
the optimal policy shifts bikes from the location with lower expected demand
to the one with higher demand, but avoids excessive movement because of the
associated cost. In the modified problem with one free bike from the first
location to the second and parking penalties above ten bikes, the optimal
policy changes: the agent is more willing to move bikes from location 1 to
location 2 (to exploit the free transfer) and tries to keep overnight bike
counts at or below ten to avoid the extra parking cost.

\vspace{1em}

\noindent\textbf{Conclusion:}\\
This experiment illustrates how MDPs and dynamic programming can be used to
solve sequential decision-making problems with stochastic dynamics. The
gridworld example makes the effect of different reward functions on the
optimal policy very clear, while the Gbike rental problem demonstrates policy
iteration in a realistic resource allocation setting. Together, these
experiments provide practical insight into value functions, Bellman equations
and the role of discounting in reinforcement learning.

\end{document}
